
A análise de requisitos é o levantamento de todas as propriedades, características e necessidades do sistema para que esse seja capaz de executar, com a qualidade desejada, as funcionalidades previstas. Essa seção serve de base tanto para o cliente, guiando as expectativas quanto ao projeto a partir da definição sólida de quais serão suas especificações, como para os projetistas, que, ao desenvolverem o sistema de acordo com os requisitos levantados, estarão seguros quanto à sua efetividade para o problema.

Os requisitos funcionais, dispostos na seção 2.2, são propriedades do sistema responsáveis por explicitar suas funcionalidades e como estas atendem às solicitações do projeto. A partir da leitura dos requisitos funcionais, tem-se noção do que o sistema deve ser capaz de executar e quais são as necessidades que ele busca suprir.

Os requisitos não-funcionais, dispostos na seção 2.3, são propriedades do sistema que mostram especificações e características relacionadas ao modo como os requisitos funcionais são operados. A partir da leitura dos requisitos não-funcionais, tem-se noção de informações técnicas e atributos de qualidade de execução das funcionalidades do sistema.

\subsection{Convenções, termos e abreviações}

\subsubsection{Identificação dos Requisitos}

Por convenção, a referência a requisitos é feita através do nome da subseção na qual eles
estão descritos, seguidos do identificador do requisito, de acordo com o seguinte modelo:
\begin{center}
    [ nome da subseção/identificador do requisito ]
\end{center}

Os requisitos serão identificados com um identificador único, com numeração iniciada em
[RF01] para um requisito funcional e [NF01] para um requisito não-funcional.

\subsubsection{Propriedades dos requisitos}

Para estabelecer a prioridade dos requisitos, nas seções 2.2 e 2.3, foram adotadas as
denominações “essencial”, “importante” e “desejável”, tal que:

\textbf{Essencial} é o requisito sem o qual o sistema não entra em funcionamento. Requisitos
essenciais são aqueles imprescindíveis, que devem ser implementados impreterivelmente.

\textbf{Importante} é o requisito sem o qual o sistema entra em funcionamento, mas de forma não
satisfatória. Requisitos importantes devem ser implementados, mas, se não forem, o sistema
poderá ser implementado e usado normalmente.

\textbf{Desejável} é o requisito que não compromete as funcionalidades básicas do sistema, isto é,
o sistema pode funcionar de forma satisfatória sem ele. Requisitos desejáveis podem ser deixados
para versões posteriores do sistema, caso não haja tempo hábil para implementação dos mesmos
nesta etapa.

\subsection{Requisitos Funcionais}

\begin{rfreq}{Requisito...}
    Descrição do Requisito...

    Prioridade: Essencial.
\end{rfreq}

\begin{rfreq}{Próximo requisito...}
    Descrição do requisito...
    
    Prioridade: Importante.
\end{rfreq}
%%
\newpage
\subsection{Requisitos Não-Funcionais}

\begin{nfreq}{Requisito não funcional}
    Descrição do requisito...

    Prioridade: Essencial.
\end{nfreq}

\begin{nfreq}{Requisito não funcional}
    Descrição do requisito...

    Prioridade: Desejável.
\end{nfreq}
